
\begin{longtable}{ | p{4.0cm} | p{8.4cm} | c | l | }
\hline\rowcolor{lightgray}
名称 & 描述 & 地址 & 访问 \\ \hline
\endfirsthead
\hline
\rowcolor{lightgray}
名称 & 描述 & 地址 & 访问 \\
\hline
\endhead
\endfoot

\multicolumn{4}{|l|}{\textbf{机器信息 CSR}} \\\hline
\hyperref[regdesc:MVENDORID]{mvendorid} & 机器供应商 ID & 0xF11 & RO \\\hline
\hyperref[regdesc:MARCHID]{marchid} & 机器架构 ID & 0xF12 & RO \\\hline
\hyperref[regdesc:MIMPID]{mimpid} & 机器实现 ID & 0xF13 & RO \\\hline
\hyperref[regdesc:MHARTID]{mhartid} & 机器 hart ID & 0xF14 & RO \\\hline
\multicolumn{4}{|l|}{\textbf{机器陷阱设置 CSR}} \\\hline
\hyperref[regdesc:MSTATUS]{mstatus} & 机器模式状态 & 0x300 & R/W \\\hline
\hyperref[regdesc:MISA]{misa} \footnote{尽管 \hyperref[regdesc:MISA]{misa} 具有读/写属性，但由于它的域是硬连线的，所以写操作无效。在 RISC-V 术语中称为 WARL（写入任意数值读取合法数值）。} & 机器 ISA & 0x301 & R/W \\\hline
\hyperref[regdesc:MIDELEG_REG]{mideleg} & 机器中断委托寄存器（在 CLIC 模式下无效） & 0x303 & RO \\\hline
\hyperref[regdesc:MIE_REG]{mie} & 机器中断使能寄存器（在 CLIC 模式下无效） & 0x304 & RO \\\hline
\hyperref[regdesc:MTVEC]{mtvec} \footnote{\hyperref[regdesc:MTVEC]{mtvec} 保存了陷阱向量配置，包括向量基地址和向量模式。} & 机器陷阱向量 & 0x305 & R/W \\\hline
\hyperref[regdesc:MCOUNTEREN]{mcounteren} & 机器计数器使能 & 0x306 & R/W \\\hline
\hyperref[regdesc:MTVT]{mtvt} & 机器向量中断基地址（参见 CLIC 规范） & 0x307 & R/W \\\hline
\multicolumn{4}{|l|}{\textbf{机器陷阱处理 CSR}} \\\hline
\hyperref[regdesc:MSCRATCH]{mscratch} & 机器暂存 & 0x340 & R/W \\\hline
\hyperref[regdesc:MEPC]{mepc} & 机器陷阱程序计数器 & 0x341 & R/W \\\hline
\hyperref[regdesc:MCAUSE]{mcause} \footnote{\hyperref[regdesc:MCAUSE]{mcause} 中反映的外部中断 ID 也包括 RISC-V 标准为核心内部中断源预留的 ID。} & 机器陷阱原因 & 0x342 & R/W \\\hline
\hyperref[regdesc:MTVAL]{mtval} & 机器陷阱值 & 0x343 & R/W \\\hline
\hyperref[regdesc:MIP_REG]{mip} & 机器中断挂起（在 CLIC 模式下无效） & 0x344 & RO \\\hline
\hyperref[regdesc:MNXTI_REG]{mnxti} & 中断处理程序地址和使能（参见 CLIC 规范） & 0x345 & R/W \\\hline
\hyperref[regdesc:MINTSTATUS]{mintstatus} & 当前中断级别（参见 CLIC 规范） & 0x346 & RO \\\hline
\hyperref[regdesc:MSCRATCHCSW_REG]{mscratchcsw} & 特权模式改变时的条件暂存交换（参见 CLIC 规范） & 0x348 & R/W \\\hline
\hyperref[regdesc:MSCRATCHCSWL_REG]{mscratchcswl} & 级别改变时的条件暂存交换（参见 CLIC 规范） & 0x349 & R/W \\\hline
\hyperref[regdesc:MCLICBASE_REG]{mclicbase} & 机器模式中断控制器基地址寄存器（自定义）（参见 CLIC 规范） & 0x350 & RO \\\hline
\multicolumn{4}{|l|}{\textbf{物理存储器保护 (PMP) CSR}} \\\hline\hyperref[sec:riscvcpu-pmp-reg-desc]{pmpcfg0} & 物理存储器保护配置 & 0x3A0 & R/W \\\hline
\hyperref[sec:riscvcpu-pmp-reg-desc]{pmpcfg1} & 物理存储器保护配置 & 0x3A1 & R/W \\\hline
\hyperref[sec:riscvcpu-pmp-reg-desc]{pmpcfg2} & 物理存储器保护配置 & 0x3A2 & R/W \\\hline
\hyperref[sec:riscvcpu-pmp-reg-desc]{pmpcfg3} & 物理存储器保护配置 & 0x3A3 & R/W \\\hline
\hyperref[sec:riscvcpu-pmp-reg-desc]{pmpaddr\regindex{n} (\regindex{n}: 0-15)} & 物理存储器保护地址寄存器 & {0x3B0+0x1*\regindex{n}} & R/W \\\hline
\multicolumn{4}{|l|}{\textbf{触发器模块 CSR（与调试模式共享）}} \\\hline
\hyperref[regdesc:TSELECT]{tselect} & 触发器选择寄存器 & 0x7A0 & R/W \\\hline
\hyperref[regdesc:TDATA1]{tdata1} & 触发器抽象数据寄存器 1 & 0x7A1 & R/W \\\hline
\hyperref[regdesc:TDATA2]{tdata2} & 触发器抽象数据寄存器 2 & 0x7A2 & R/W \\\hline
\hyperref[regdesc:TDATA3]{tdata3} & 触发器抽象数据寄存器 3 & 0x7A3 & R/W \\\hline
\hyperref[regdesc:TINFO]{tinfo} & 触发器信息寄存器 & 0x7A4 & R/W \\\hline
\hyperref[regdesc:TCONTROL]{tcontrol} & 全局触发器控制 & 0x7A5 & R/W \\\hline
\hyperref[regdesc:MCONTEXT]{mcontext} & 机器模式上下文寄存器 & 0x7A8 & R/W \\\hline
\multicolumn{4}{|l|}{\textbf{调试模式 CSR（仅在调试模式下可访问）}} \\\hline
\hyperref[regdesc:DCSR]{dcsr} & 调试控制和状态 & 0x7B0 & R/W \\\hline
\hyperref[regdesc:DPC]{dpc} & 调试 PC & 0x7B1 & R/W \\\hline
\hyperref[regdesc:DSCRATCH0]{dscratch0} & 调试暂存寄存器 0 & 0x7B2 & R/W \\\hline
\hyperref[regdesc:DSCRATCH1]{dscratch1} & 调试暂存寄存器 1 & 0x7B3 & R/W \\\hline
\hyperref[regdesc:dmcs2]{dmcs2} & 调试模式状态控制寄存器 2 & 0x32 & R/W \\\hline
\multicolumn{4}{|l|}{\textbf{机器计数器设置}} \\\hline
\hyperref[regdesc:MCOUNTINHIBIT]{mcountinhibit} & 机器计数器控制寄存器 & 0x320 & R/W \\\hline
\hyperref[regdesc:MHPMEVENT8]{mhpmevent8} & 机器性能监控事件选择器 & 0x308 & R/W \\\hline
\hyperref[regdesc:MHPMEVENT9]{mhpmevent9} & 机器性能监控事件选择器 & 0x329& R/W \\\hline
\hyperref[regdesc:MHPMEVENT13]{mhpmevent13} & 机器性能监控事件选择器 & 0x32D & R/W \\\hline
\multicolumn{4}{|l|}{\textbf{机器计数器/定时器}} \\\hline
\hyperref[regdesc:MCYCLE]{mcycle} & 机器周期计数器 & 0xB00 & R/W \\\hline
\hyperref[regdesc:MINSTRET]{minstret} & 机器指令退役计数器 & 0xB02 & R/W \\\hline
\hyperref[regdesc:MHPMCOUNTER8]{mhpmcounter8} & 机器性能监控计数器 & 0xB08 & R/W \\\hline
\hyperref[regdesc:MHPMCOUNTER9]{mhpmcounter9} & 机器性能监控计数器 & 0xB09 & R/W \\\hline
\hyperref[regdesc:MHPMCOUNTER13]{mhpmcounter13} & 机器性能监控计数器 & 0xB0D & R/W \\\hline
\hyperref[regdesc:MCYCLEH]{mcycleh} & \hyperref[regdesc:MCYCLE]{mcycle} 的高 32 位 & 0xB80 & R/W \\\hline
\hyperref[regdesc:MINSTRETH]{minstreth} & \hyperref[regdesc:MINSTRET]{minstret} 的高 32 位 & 0xB82 & R/W \\\hline
\hyperref[regdesc:MHPMCOUNTER8H]{mhpmcounter8h} & \hyperref[regdesc:MHPMCOUNTER8]{mhpmcounter8} 的高 32 位 & 0xB88 & R/W \\\hline
\hyperref[regdesc:MHPMCOUNTER9H]{mhpmcounter9h} & \hyperref[regdesc:MHPMCOUNTER9]{mhpmcounter9} 的高 32 位 & 0xB89 & R/W \\\hline
\hyperref[regdesc:MHPMCOUNTER13H]{mhpmcounter13h} & \hyperref[regdesc:MHPMCOUNTER13]{mhpmcounter13} 的高 32 位 & 0xB8D & R/W \\\hline
\multicolumn{4}{|l|}{\textbf{非特权/用户模式浮点 CSR}} \\\hline
\hyperref[regdesc:FFLAGS]{fflags} & 浮点累积异常 & 0x001 & R/W \\\hline
\hyperref[regdesc:FRM]{frm} & 浮点动态舍入模式 & 0x002 & R/W \\\hline
\hyperref[regdesc:FCSR]{fcsr} & 浮点控制和状态寄存器 (frm + fflags) & 0x003 & R/W \\\hline
\multicolumn{4}{|l|}{\textbf{非特权/用户模式计数器/定时器}} \\\hline
\hyperref[regdesc:CYCLE]{cycle} & RDCYCLE 指令的周期计数器 & 0xC00 & RO\\\hline
\hyperref[regdesc:TIME]{time} & RDTIME 指令的定时器 & 0xC01 & RO \\\hline
\hyperref[regdesc:INSTRET]{instret} & RDINSTRET 指令的退役计数器 & 0xC02 & RO \\\hline
\hyperref[regdesc:HPMCOUNTER8]{hpmcounter8} & 性能监控计数器 & 0xC08 & RO \\\hline
\hyperref[regdesc:HPMCOUNTER9]{hpmcounter9} & 性能监控计数器 & 0xC09 & RO \\\hline
\hyperref[regdesc:HPMCOUNTER13]{hpmcounter13} & 性能监控计数器 & 0xC0D & RO \\\hline
\hyperref[regdesc:CYCLEH]{cycleh} & \hyperref[regdesc:CYCLE]{cycle} 的高 32 位 & 0xC80 & RO \\\hline
\hyperref[regdesc:INSTRETH]{instreth} & \hyperref[regdesc:INSTRET]{instret} 的高 32 位 & 0xC82 & RO \\\hline
\hyperref[regdesc:HPMCOUNTER8H]{hpmcounter8h} & \hyperref[regdesc:HPMCOUNTER8]{hpmcounter8} 的高 32 位 & 0xC88 & RO \\\hline
\hyperref[regdesc:HPMCOUNTER9H]{hpmcounter9h} & \hyperref[regdesc:HPMCOUNTER9]{hpmcounter9} 的高 32 位 & 0xC89 & RO \\\hline
\hyperref[regdesc:HPMCOUNTER13H]{hpmcounter13h} & \hyperref[regdesc:HPMCOUNTER13]{hpmcounter13} 的高 32 位 & 0xC8D & RO \\\hline
\multicolumn{4}{|l|}{\textbf{自定义系统 CSR（自定义）}} \\\hline
\multicolumn{4}{|l|}{\textbf{浮点用户模式 CSR（自定义）}} \\\hline
\hyperref[regdesc:FXCR]{fxcr} & 浮点扩展控制寄存器 & 0x800 & R/W \\\hline
\multicolumn{4}{|l|}{\textbf{GPIO 访问 CSR（自定义）}} \\\hline
\hyperref[regdesc:GPIO_OEN]{cpu\_gpio\_oen} & GPIO 输出使能 & 0x803 & R/W \\\hline
\hyperref[regdesc:GPIO_IN]{cpu\_gpio\_in} & GPIO 输入值 & 0x804 & RO \\\hline
\hyperref[regdesc:GPIO_OUT]{cpu\_gpio\_out} & GPIO 输出值 & 0x805 & R/W \\\hline
\multicolumn{4}{|l|}{\textbf{扩展控制和状态 CSR（自定义）}} \\\hline
\hyperref[sec:riscvcpu-hwloop-reg-desc]{mext\_ill\_reg} & 机器扩展非法寄存器 & 0x7F0 & R/W \\\hline
\hyperref[sec:riscvcpu-hwloop-reg-desc]{mhwloop\_state\_reg} & 机器模式 HWLP 扩展状态寄存器 & 0x7F1 & R/W \\\hline
\hyperref[regdesc:MEXT_PIE_STATUS]{mext\_pie\_status} & 机器模式 PIE 扩展状态寄存器 & 0x7F2 & R/W \\\hline
\multicolumn{4}{|l|}{\textbf{Zc 扩展 CSR}} \\\hline
\hyperref[regdesc:JVT]{jvt} & 机器向量跳转表基地址和控制寄存器 & 0x017 & R/W \\\hline
\multicolumn{4}{|l|}{\textbf{机器扩展状态 CSR（自定义）}} \\\hline
\hyperref[regdesc:MEXSTATUS]{mexstatus} & 机器扩展控制和状态寄存器 & 0x7E1 & R/W \\\hline
\hyperref[regdesc:MHINT]{mhint} & 机器隐式操作寄存器 & 0x7C5 & R/W \\\hline
\multicolumn{4}{|l|}{\textbf{物理存储器属性检查器 (PMAC) CSR（自定义）}} \\\hline
\hyperref[sec:riscvcpu-pma-reg-desc]{pma\_cfg\regindex{n} (\regindex{n}: 0-15)} & 物理存储器属性配置寄存器 & {0xBC0+0x1*\regindex{n}} & R/W \\\hline
\hyperref[sec:riscvcpu-pma-reg-desc]{pma\_addr\regindex{n} (\regindex{n}: 0-15)} & 物理存储器属性地址寄存器 & {0xBD0+0x1*\regindex{n}} & R/W \\\hline
\multicolumn{4}{|l|}{\textbf{机器硬件循环 CSR（自定义）}} \\\hline
\hyperref[sec:riscvcpu-hwloop-reg-desc]{mhloop0\_start\_addr} & 循环 0 的 32 位起始地址 & 0x7C6 & R/W \\\hline
\hyperref[sec:riscvcpu-hwloop-reg-desc]{mhloop0\_end\_addr} & 循环 0 的 32 位结束地址 & 0x7C7 & R/W \\\hline
\hyperref[sec:riscvcpu-hwloop-reg-desc]{mhloop0\_count} & 硬件循环 0 的循环计数 & 0x7C8 & R/W \\\hline
\hyperref[sec:riscvcpu-hwloop-reg-desc]{mhloop1\_start\_addr} & 循环 1 的 32 位起始地址 & 0x7C9 & R/W \\\hline
\hyperref[sec:riscvcpu-hwloop-reg-desc]{mhloop1\_end\_addr} & 循环 1 的 32 位结束地址 & 0x7CA & R/W \\\hline
\hyperref[sec:riscvcpu-hwloop-reg-desc]{mhloop1\_count} & 硬件循环 1 的循环计数 & 0x7CB & R/W \\\hline
\multicolumn{4}{|l|}{\textbf{用户硬件循环 CSR（自定义）}} \\\hline
\hyperref[sec:riscvcpu-hwloop-reg-desc]{uhloop0\_start\_addr} & 循环 0 的 32 位起始地址 & 0x8C0 & R/W \\\hline
\hyperref[sec:riscvcpu-hwloop-reg-desc]{uhloop0\_end\_addr} & 循环 0 的 32 位结束地址 & 0x8C1 & R/W \\\hline
\hyperref[sec:riscvcpu-hwloop-reg-desc]{uhloop0\_count} & 硬件循环 0 的循环计数 & 0x8C2 & R/W \\\hline
\hyperref[sec:riscvcpu-hwloop-reg-desc]{uhloop1\_start\_addr} & 循环 1 的 32 位起始地址 & 0x8C3 & R/W \\\hline
\hyperref[sec:riscvcpu-hwloop-reg-desc]{uhloop1\_end\_addr} & 循环 1 的 32 位结束地址 & 0x8C4 & R/W \\\hline
\hyperref[sec:riscvcpu-hwloop-reg-desc]{uhloop1\_count} & 硬件循环 1 的循环计数 & 0x8C5 & R/W \\\hline
\multicolumn{4}{|l|}{\textbf{总线错误异常上下文 CSR（自定义）}} \\\hline
\hyperref[regdesc:LDPC0]{ldpc0} & 加载总线错误 PC 寄存器 0 & 0xBE0 & R/W \\\hline
\hyperref[regdesc:LDPC1]{ldpc1} & 加载总线错误 PC 寄存器 1 & 0xBE1 & R/W \\\hline
\hyperref[regdesc:LDTVAL0]{ldtval0} & 加载总线错误访问地址寄存器 0 & 0xBE8 & R/W \\\hline
\hyperref[regdesc:LDTVAL1]{ldtval1} & 加载总线错误访问地址寄存器 1 & 0xBE9 & R/W \\\hline
\hyperref[regdesc:STPC0]{stpc0} & 存储总线错误 PC 寄存器 0 & 0xBF0 & R/W \\\hline
\hyperref[regdesc:STPC1]{stpc1} & 存储总线错误 PC 寄存器 1 & 0xBF1 & R/W \\\hline
\hyperref[regdesc:STPC2]{stpc2} & 存储总线错误 PC 寄存器 2 & 0xBF2 & R/W \\\hline
\hyperref[regdesc:STTVAL0]{sttval0} & 存储总线错误访问地址寄存器 0 & 0xBF8 & R/W \\\hline
\hyperref[regdesc:STTVAL1]{sttval1} & 存储总线错误访问地址寄存器 1 & 0xBF9 & R/W \\\hline
\hyperref[regdesc:STTVAL2]{sttval2} & 存储总线错误访问地址寄存器 2 & 0xBFA & R/W \\\hline
\end{longtable}
